\section{Data and methods}\label{sec:methods}

\subsection{Data}
\paragraph{Citation graph} The graph is the January 2024 OpenCitations snapshot distributed as parsed tables by
\citet{arjomandi2024dataset} \citep[see also][]{peroni2020opencitations}. Publication years, identifiers and journals
come from OpenCitations Meta \citep{massari2024meta,ocmeta2024dump}. Measured by the monthly count of citing papers,
the snapshot is effectively complete to about September 2023. A paper is \emph{biomedical} if it has a PMID in Meta or
its DOI appears in the PubMed 2026 baseline, from which publication types and comment/correction links are also taken.

\paragraph{Retractions} Retraction labels come from the Retraction Watch database distributed by Crossref
\citep{rittman2025rw}, in the repository snapshot of 30 October 2026 (62,735 retracted papers). Papers were linked to graph
nodes by DOI or PMID; 52,254 of the 62,735 (83.3\%) were matched, giving 53,170 graph nodes. The linked table is
public \citep{arjomandi2026retracteddata}. The date of the first retraction notice and the union of the reasons given
by Retraction Watch are kept per paper; reinstated papers are excluded.

\paragraph{Reason classes} Each retracted paper receives one class, assigned in this order: \emph{paper mill} (reason
``Paper Mill''); \emph{misconduct} (fabrication or falsification, manipulation of images, data or results, plagiarism,
duplicated images, compromised peer review, rogue editor, forged authorship or affiliation, computer-generated
content, hoax papers, or an official misconduct finding); \emph{honest error} (any ``Error in'' or ``Error by''
reason, results not reproducible, contaminated materials or cell lines, ``Doing the Right Thing''); and \emph{other}
(everything else, including notices with limited information). The full list of reason strings is in the
specification file (see the data and code availability statement).

\subsection{Populations}
\paragraph{Eligible papers} Biomedical papers published 2000--2017 with a journal in Meta, at least one reference in
the graph, and no duplicate DOI record. The year limit ensures that the five-year window ends by 2022, inside the
complete part of the snapshot. Papers without references are excluded: their index equals 1 whenever they are cited
and is undefined otherwise, and they are the records behind the CD$_5 = 1$ artefact discussed by
\citet{holst2024artefacts}. Of 12,522,210 biomedical papers from 2000--2017 with a journal, 341,432 (2.7\%) have no
reference in the graph, leaving 12,180,778 eligible papers.

\paragraph{Retracted papers} Eligible papers matched to a retraction: 10,450. The \emph{primary set} contains the
4,587 papers whose first retraction notice is dated after 31 December of the fifth year after publication, so that the
whole scoring window precedes the retraction. The 5,863 papers retracted earlier form, together with the primary set,
the \emph{secondary set} used in a robustness check.

\paragraph{Controls} For every retracted paper, five eligible papers from the same journal and publication year that
are neither in Retraction Watch nor flagged as retracted in PubMed. Within a stratum, controls are the papers with the
lowest value of a fixed hash of their identifier (seed 20261109), which makes the draw deterministic. This gives
52,147 controls in 5,458 journal $\times$ year strata; four strata (four retracted papers) have no eligible control.
This pool is drawn for all 10,450 retracted papers, not just the 4,587 primary-set papers. A comparison uses only the
controls with a defined score in strata containing scored papers of the compared set. The primary analysis uses
2,239 strata containing 4,586 primary-set papers and 31,662 scored controls.

\paragraph{Reference sample} A deterministic 0.2\% sample of all eligible papers (24,491 papers) provides the
yearly distribution of the index, from which the top-5\% and top-1\% thresholds are taken.

\paragraph{Post hoc population} The same rules applied to papers published 2000--2019 with a three-year window give
13,661 retracted papers (7,350 retracted after their window), 68,190 controls and a reference sample of 28,784 papers.
It contains every paper of the population above.

\subsection{The CD index}\label{sec:score}
For a focal paper $f$ published in year $t$ with references $R(f)$ in the graph, let $A$ be the papers published in the
window that cite $f$, and $B$ the papers published in the window, other than $f$, that cite at least one paper in
$R(f)$. With $n_i = |A \setminus B|$ (citing papers that cite none of the references), $n_j = |A \cap B|$ (citing
papers that also cite at least one reference) and $n_k = |B \setminus A|$ (papers that cite only the references),
\begin{equation}
\mathrm{CD}(f) = \frac{n_i - n_j}{n_i + n_j + n_k},
\end{equation}
undefined when the denominator is zero \citep{funk2017dynamic,park2023papers}. Self-loops are ignored.
\citet{park2023papers} measure the index five years after each paper's publication year. The main window (CD$_5$)
covers the citing years $t$ to $t+5$; because year-level data cannot order the focal paper and other papers of the same
year, those papers are included. The reference implementation of the index counts only citing works with a later
timestamp than the focal work\footnote{Software \texttt{cdindex} by R.~J. Funk,
\url{https://github.com/russellfunk/cdindex}, file \texttt{src/cdindex.c} at commit \texttt{4eece64}.}, which with
year-level timestamps corresponds to the citing years $t+1$ to $t+5$; this window is used in robustness check (d).

The post hoc variant without $n_k$,
\begin{equation}
\mathrm{DI}_{nok}(f) = \frac{n_i - n_j}{n_i + n_j},
\end{equation}
was proposed by \citet{wuyan2019solo} and studied as DI$_1^{nok}$ by \citet{bornmann2020convergent}. It depends only on
the papers that cite $f$ and is defined when there is at least one; $(1 - \mathrm{DI}_{nok})/2$ is the share of
citing papers that also cite a reference. CD and DI$_{nok}$ always have the same sign.

The score script reads only the list of focal papers and their years; it never reads retraction labels. Its output
(86,950 papers) was hashed (SHA-256) before the evaluation script joined the labels, and the evaluation script checks
the hash. An implementation in DuckDB \citep{raasveldt2019duckdb} was tested against a brute-force count on random
graphs (1,200 focal papers, both windows, identical counts). For the before-and-after comparison, a second label-free
script counted the citing papers of each focal paper by citing year; its sums over the scoring windows equal the frozen
counts for all 110,462 papers of both populations, and its output was hashed before labels were joined.

\paragraph{Pre-specification record} The unchanged specification is \texttt{SPEC.md}, committed in the then-private
project repository on 10 November 2026 at 04:02:19 UTC (Git commit
\texttt{293755b1cb7087d6f226d6f96f4b97cb0148e4c5}). Its SHA-256,
\nolinkurl{1403d0d13d35bd32c40ede1be24833ab15e7977edb2f56c2128ef8823373ad31},
is recorded in the score run's \texttt{run.json}, which reports a start time of 06:12:02 UTC the same day
(\texttt{e03\_cd5\_20261110\_01}); the evaluation run records the same hash. These files and provenance are available
with the code. Git dates and local run records are not an independent registry timestamp, so the design is described
as pre-specified, not publicly pre-registered.

\subsection{Analyses}
All analyses below were specified before any score was computed, except those marked post hoc. The plan of each post
hoc analysis, with its margins and expected directions, was recorded before it was run.
\begin{itemize}
\item \textbf{RQ1 (primary):} the mean difference in CD$_5$ between retracted papers and their controls within
journal $\times$ year strata, averaged over retracted papers, with a 95\% interval from 2,000 bootstrap resamples of
strata; and the Mantel--Haenszel odds ratio \citep{mantel1959} for a positive score (CD$_5 > 0$), with the
Robins--Breslow--Greenland variance \citep{robins1986estimator}.
\item \textbf{RQ2:} the same by reason class.
\item \textbf{RQ3:} the share of papers at or above the 95th and 99th percentiles of CD$_5$ in their publication year
(from the reference sample), with Wilson intervals \citep{wilson1927}, and the stratified odds ratio for reaching them.
\item \textbf{Robustness (pre-specified):} (a) without CD$_5 = 1$ papers; (b) without reviews (PubMed publication
types Review, Systematic Review, Meta-Analysis); (c) strata further split by deciles of reference count; (d) window
$t+1$ to $t+5$; (e) the secondary set; (f) papers with at least five citations in the window.
\item \textbf{Post hoc, first round:} (P1) reference and citation counts by group; (P2) a three-year window (CD$_3$)
for papers published 2000--2019, which adds paper-mill papers from 2018--2019, with the primary set defined by a
retraction after the third year.
\item \textbf{Post hoc, second round:} (P3) equivalence tests: two one-sided tests \citep{schuirmann1987tost,
lakens2017equivalence} with margins of $\pm 0.1$ standard deviations of the score in the reference sample for mean
differences and 0.80--1.25 for odds ratios; these post hoc margins are conventions, not validated thresholds for the CD
index. Equivalence is shown if the 90\% interval lies inside the margin; the smallest supported symmetric margin is
reported so readers can apply a different threshold. (P4) A positive control: reviews compared with the other
never-retracted controls of the same journal and year, expected to be more consolidating because papers citing a
review often also cite the work it reviews. (P5) RQ1--RQ3 with DI$_{nok}$. (P6) Before and after the notice: for
retracted papers of the 2000--2019 population with retraction year $R$ between $t+1$ and 2021, the change in
DI$_{nok}$ from the citing years $R-3$ to $R-1$ (not before $t$) to the years $R+1$ to $R+3$ (not after 2022),
minus the mean change of their stratum's controls over the same years, for papers with at least one citing paper in
both periods, with year $R$ and notice records excluded.
(P7) The sign shift: papers without citing papers in the window, odds ratios among cited papers only, and conditional
logistic regression within journal $\times$ year strata with $\log(1+\text{references})$ and
$\log(1+\text{citing papers in the window})$ as covariates. (P8) Editorial notices: citing records that PubMed
types as retraction notices, published errata or expressions of concern, or that carry the corresponding links
(RetractionOf, ErratumFor, ExpressionOfConcernFor), and robustness check (e) without them.
\item \textbf{Post hoc, third round:} (P9) checks of P6 for an earlier trend: a placebo notice three years before the
real one, and yearly means on a balanced set of papers; (P10) the size of $n_k$.
\item \textbf{Post hoc, fourth round (descriptive):} (P11) the controls behind the primary comparisons: their number in
the strata used, and their shares above zero and above the yearly thresholds, with each stratum weighted by its
retracted papers as for the mean difference (95\% interval from 2,000 bootstrap resamples of strata).
\end{itemize}
Only group-level results are reported. No per-paper score of a non-retracted paper is released.

\paragraph{Software and AI assistance} The analyses ran in Python 3.11 with DuckDB 1.5, pandas 3.1, NumPy 2.4, SciPy
1.17 and statsmodels 0.15 on one cloud machine; figures were drawn with Matplotlib 3.10 from the run records. The
analysis and figure scripts were written with the help of a generative AI assistant (Claude Opus 5.5, Anthropic) under
the author's direction. Every counting step was tested against brute-force counts on random graphs, each evaluation
script first reproduces the earlier estimates it builds on, and all scripts and run records are public
(see the data and code availability statement).
